% ════════════════════════════════════════════
% 测算方法（5.4.3）
% ════════════════════════════════════════════
\subsubsection{测算方法}

本节从预测还原、效应定义、区间构造与检验四个方面建立效应测算框架，把回归系数翻译为曝光次数、倍数与百分比三种口径。

最终模型的因变量为 log1p 曝光量，因此在给定协变量后，线性预测值只能解释为对数尺度上的条件均值。若直接取指数，得到的是条件几何均值，会系统性低估曝光次数的算术均值。Duan 提出的烟雾还原用残差的经验分布修正这一偏差\cite{duan1983}：先对每条训练样本计算残差，再取残差指数的平均值作为烟雾因子，本样本的烟雾因子为 2.1896，明显大于 1；把线性预测值取指数后乘以烟雾因子再减 1，即得到曝光次数的条件均值估计。该做法不需要假定残差服从正态分布。
\begin{equation}
    \hat m(\boldsymbol{x})=e^{\hat\eta(\boldsymbol{x})}\hat\lambda-1,
    \qquad
    \hat\lambda=\frac{1}{n}\sum_{i=1}^{n}e^{e_i},
\end{equation}
其中，$\hat\eta(\boldsymbol{x})$ 为对数尺度预测值，$e_i$ 为训练样本残差，$\hat\lambda$ 为烟雾因子。

对二元处理，第 $i$ 条帖子的个体处理效应为两种情景预测值之差，样本平均处理效应为 500 条帖子个体效应的平均。相应的倍数用处理组预测均值之和除以对照组预测均值之和，百分比则用处理系数的指数减一乘以 100。前者是还原到曝光尺度后的均值比，后者是控制其他变量后的净百分比，两者数值接近但不完全相同，正文同时报告以避免口径混淆。

区间估计采用参数模拟。从小样本渐近理论出发，系数估计近似服从以点估计为均值、以 HC3 协方差为协方差矩阵的多元正态分布。独立抽取 1000 组系数，每组计算代表性画像与全样本的效应，取 2.5\% 与 97.5\% 分位数作为 95\% 置信区间。与仅用系数标准误的区间相比，该方法同时传播了多个系数的协方差，适合倍数与差值这类非线性函数。

内容主题的联合检验以“两个主题哑变量同时为零”为原假设。构造约束矩阵后，Wald 统计量由约束向量与协方差矩阵的二次型给出，服从自由度为 2 的卡方分布；若统计量超过临界值，则拒绝主题整体无影响的判断。两两比较主题 2 与主题 3 时，约束为两个系数之差为零，用系数协方差计算差值标准误，构造 z 统计量。

为判断结论是否依赖模型设定，在六种设定下重复计算付费与视频效应：M1 全模型、M2 月份固定效应、M3 仅粉丝规模对数、M4 粉丝规模对数加线性月份、Huber 稳健回归与 Gamma 广义线性模型。前四种改变解释变量集合，后两种改变估计方法。若六种设定的效应方向一致、区间均不跨过 0 或 1，则认为结论稳健。

综上所述，本节建立了由烟雾还原、平均处理效应、参数模拟与 Wald 检验组成的效应测算框架，得到了可解释为曝光次数与倍数的输出口径，接下来将基于此给出具体结果。
